% arXiv-ready LaTeX paper: Rank-based Optional Stopping via Structure Towers
% Compiled with pdflatex (no shell-escape, no minted)

\documentclass[11pt]{article}

% Minimal package set
\usepackage{amsmath,amsthm,amssymb,mathtools}
\usepackage[colorlinks=true,linkcolor=blue,citecolor=blue,urlcolor=blue]{hyperref}
\usepackage[margin=1in]{geometry}

% Theorem environments
\newtheorem{theorem}{Theorem}[section]
\newtheorem{lemma}[theorem]{Lemma}
\newtheorem{proposition}[theorem]{Proposition}
\newtheorem{corollary}[theorem]{Corollary}
\theoremstyle{definition}
\newtheorem{definition}[theorem]{Definition}
\newtheorem{example}[theorem]{Example}
\theoremstyle{remark}
\newtheorem{remark}[theorem]{Remark}

% Macros for key objects
\newcommand{\layer}{\mathcal{L}}
\newcommand{\minLayer}{\mathrm{minLayer}}
\newcommand{\rank}{\mathrm{rank}}
\newcommand{\Tower}{\mathsf{Tower}}
\newcommand{\Filtration}{\mathcal{F}}
\newcommand{\NN}{\mathbb{N}}
\newcommand{\QQ}{\mathbb{Q}}
\newcommand{\RR}{\mathbb{R}}
\newcommand{\EE}{\mathbb{E}}
\newcommand{\PP}{\mathbb{P}}
\newcommand{\WithTop}[1]{#1 \cup \{\top\}}

% Title and authors
\title{Rank-based Optional Stopping via Structure Towers}

\author{%
  Generated Collaboratively\\
  \small\textit{TeX generation: Claude Code (Anthropic)}\\
  \small\textit{Lean formalization: Codex (OpenAI)}%
}

\date{\today}

\begin{document}

\maketitle

\begin{abstract}
We present a formalization of the Optional Stopping Theorem (OST) for discrete-time martingales using \emph{structure towers} and \emph{rank functions}.
Structure towers provide a layer-based view of filtered probability spaces, where each layer corresponds to a stopping set.
We establish that stopping times are canonically identified with rank functions via the correspondence $\tau(\omega) = \minLayer(\omega)$.
This identification allows a purely algebraic proof of the bounded OST.
The theory is fully mechanized in Lean~4 using Mathlib, yielding executable decision procedures for finite probability spaces.
We isolate category-theoretic aspects in an appendix, keeping the main exposition accessible to probability theorists and formal-methods practitioners.
\end{abstract}

\section{Introduction}
\label{sec:intro}

The Optional Stopping Theorem is a cornerstone of martingale theory, asserting that under suitable conditions, the expectation of a martingale at a stopping time equals its initial expectation.
Classical proofs rely on measure-theoretic arguments and filtration adaptedness.
We propose an alternative framework based on \emph{structure towers} and \emph{rank functions}, inspired by Bourbaki's mother-structure philosophy.
Our main contributions are:

\begin{itemize}
\item \textbf{Rank--tower correspondence} (Section~\ref{sec:tower}): A bijection between structure towers with minimal layers and rank functions $\rho: X \to \NN$.
\item \textbf{Stopping times as ranks} (Section~\ref{sec:stopping-rank}): A translation showing that a stopping time $\tau: \Omega \to \NN$ is a rank function on the probability space, with stopping sets as tower layers.
\item \textbf{Rank OST} (Section~\ref{sec:rank-ost}): A proof of the bounded Optional Stopping Theorem formulated entirely in rank-theoretic language, with layer separation from implementation details.
\item \textbf{Lean mechanization} (Section~\ref{sec:lean}): A complete formalization in Lean~4, including executable verification for finite spaces with rational probabilities.
\end{itemize}

The rank-based view unifies concepts from algebra (layers, monotonicity), order theory (minimal elements), and probability (stopping times, martingales).
We restrict to \emph{discrete settings}: finite sample spaces $\Omega$, rational probabilities, and bounded stopping times, ensuring computability and avoiding technical measure-theoretic overhead.
Category-theoretic structures (functors, adjunctions) are deferred to Appendix~\ref{app:cat}.

\paragraph{Related work.}
Formalization of martingale theory has been pursued in Isabelle/HOL~\cite{isabelle-ost} and Coq~\cite{coq-prob}.
Our rank-based approach is novel in treating stopping times as \emph{rank functions}, enabling a compositional, layer-wise proof strategy.

\paragraph{Notation.}
Once introduced, notation is fixed throughout.
$\NN = \{0,1,2,\ldots\}$; $\WithTop{\NN} = \NN \cup \{\top\}$ with $n < \top$ for all $n \in \NN$.

\section{Discrete Setting}
\label{sec:discrete}

We work in a simplified probability framework suitable for mechanization and computation.
Fix a finite set $\Omega$ (the sample space), a $\sigma$-algebra $\mathcal{A}$ over $\Omega$, and a probability measure $\mu: \mathcal{A} \to \QQ_{\geq 0}$ with $\mu(\Omega) = 1$.
A \emph{filtration} is an increasing sequence $\Filtration = (\Filtration_n)_{n \in \NN}$ of sub-$\sigma$-algebras: $\Filtration_0 \subseteq \Filtration_1 \subseteq \cdots \subseteq \mathcal{A}$.
A stochastic process $X = (X_n)_{n \in \NN}$ is \emph{adapted} if $X_n$ is $\Filtration_n$-measurable for each $n$.

\begin{definition}[Martingale]
\label{def:martingale}
An adapted, integrable process $M = (M_n)_{n \in \NN}$ is a \emph{martingale} if
\[
\EE[M_{n+1} \mid \Filtration_n] = M_n \quad \text{for all } n \in \NN.
\]
\end{definition}

\begin{definition}[Stopping time]
\label{def:stopping-time}
A function $\tau: \Omega \to \NN$ is a \emph{stopping time} (relative to $\Filtration$) if
\[
\{\omega \in \Omega : \tau(\omega) \leq n\} \in \Filtration_n \quad \text{for all } n \in \NN.
\]
\end{definition}

The set $\{\tau \leq n\}$ is called the \emph{$n$-th stopping set}.
Intuitively, $\tau(\omega)$ is the first time the ``event $\omega$'' becomes observable.

\begin{definition}[Stopped process]
\label{def:stopped-process}
For a process $X$ and stopping time $\tau$, the \emph{stopped process} $X^\tau$ is
\[
X^\tau_n(\omega) = X_{\min(n, \tau(\omega))}(\omega).
\]
\end{definition}

Under finiteness assumptions, $X^\tau$ is adapted to $\Filtration$ and inherits integrability.

\section{Minimal Structure Towers}
\label{sec:tower}

Structure towers formalize the idea of ``layered sets'' with a canonical minimal-layer function.
They provide the algebraic skeleton for our rank-theoretic framework.

\begin{definition}[Structure tower with minimal layer]
\label{def:tower}
A \emph{structure tower} $\Tower$ consists of:
\begin{itemize}
\item A carrier set $X$,
\item A family of subsets $\layer_n \subseteq X$ for $n \in \NN$,
\item \textbf{Covering}: $\forall x \in X, \, \exists n \in \NN, \, x \in \layer_n$,
\item \textbf{Monotonicity}: $\layer_i \subseteq \layer_j$ whenever $i \leq j$,
\item A function $\minLayer: X \to \NN$ such that:
  \begin{enumerate}
  \item $\forall x \in X, \, x \in \layer_{\minLayer(x)}$,
  \item $\forall x \in X, \, \forall n \in \NN, \, x \in \layer_n \Rightarrow \minLayer(x) \leq n$.
  \end{enumerate}
\end{itemize}
\end{definition}

In Lean, this is encoded as a \texttt{structure StructureTowerWithMin}.
The $\minLayer$ function captures the ``first layer'' in which $x$ appears.

\subsection{Rank functions and towers: the canonical correspondence}

A rank function assigns a natural-number rank to each element, inducing a tower via sublevel sets.

\begin{definition}[Tower from rank]
\label{def:tower-from-rank}
Given $\rho: X \to \WithTop{\NN}$ with $\forall x, \, \exists n \in \NN, \, \rho(x) \leq n$, define $\Tower_\rho$ by
\[
\layer_n = \{x \in X : \rho(x) \leq n\}, \quad \minLayer(x) = \min\{n \in \NN : \rho(x) \leq n\}.
\]
\end{definition}

\begin{lemma}[Characterization of layers]
\label{lem:tower-from-rank-layer}
$x \in \layer_n \iff \rho(x) \leq n$.
\end{lemma}

\begin{lemma}[minLayer equals rank]
\label{lem:minlayer-eq-rank}
If $\rho: X \to \NN$ (taking values in $\NN$, not $\WithTop{\NN}$), then $\minLayer_{\Tower_\rho}(x) = \rho(x)$.
\end{lemma}

\noindent\textbf{Proof sketch.}
By definition, $\minLayer(x)$ is the least $n$ with $\rho(x) \leq n$.
Since $\rho(x) \in \NN$, $n = \rho(x)$ is the unique minimum.
\qed

Conversely, any tower induces a rank function.

\begin{definition}[Rank from tower]
\label{def:rank-from-tower}
Given $\Tower$, define $\rank_\Tower: X \to \WithTop{\NN}$ by $\rank_\Tower(x) = \minLayer(x)$.
\end{definition}

\begin{theorem}[Rank--tower isomorphism]
\label{thm:rank-tower-iso}
The constructions $\rho \mapsto \Tower_\rho$ and $\Tower \mapsto \rank_\Tower$ are mutually inverse (up to isomorphism):
\begin{enumerate}
\item For $\rho: X \to \NN$, $\rank_{\Tower_\rho} = \rho$ pointwise.
\item For $\Tower$, $\layer_n^{\Tower_{\rank_\Tower}} = \layer_n$ for all $n$.
\end{enumerate}
\end{theorem}

\noindent\textbf{Proof sketch.}
(1) follows from Lemma~\ref{lem:minlayer-eq-rank}.
(2): $x \in \layer_n^{\Tower_{\rank_\Tower}}$ iff $\rank_\Tower(x) \leq n$ iff $\minLayer(x) \leq n$ iff $x \in \layer_n$ by minimality.
\qed

This bijection justifies identifying towers with rank functions.
In Lean, this is:

\texttt{noncomputable def towerFromRank (rho : X -> WithTop Nat) ...}

\begin{example}[Linear span rank for $\QQ^2$]
\label{ex:vec2q}
Let $X = \QQ^2$ and define $\rho(v)$ as the minimum number of basis vectors needed to represent $v$:
\[
\rho(0,0) = 0, \quad \rho(v) = 1 \text{ if $v$ is on an axis}, \quad \rho(v) = 2 \text{ otherwise}.
\]
Then $\layer_0 = \{(0,0)\}$, $\layer_1 = \{v : v_1 = 0 \text{ or } v_2 = 0\}$, $\layer_2 = \QQ^2$.
\end{example}

The correspondence provides a \emph{normal form}: every tower is the tower-from-rank of its $\minLayer$ function.

\section{Stopping Times as Ranks}
\label{sec:stopping-rank}

We now translate stopping times into rank functions, making explicit the layer-based structure of filtrations.
Let $(\Omega, \Filtration, \mu)$ be a filtered finite probability space.

\begin{definition}[Stopping time as rank function]
\label{def:stopping-as-rank}
For a stopping time $\tau: \Omega \to \NN$, define $\rank_\tau: \Omega \to \NN$ by $\rank_\tau(\omega) = \tau(\omega)$.
\end{definition}

Trivial as it seems, this definition is conceptually significant: it treats $\tau$ not as a random variable but as a \emph{rank function} on $\Omega$.

\begin{proposition}[Tower from stopping time]
\label{prop:tower-from-stopping-time}
Given a stopping time $\tau$, the structure $\Tower_\tau$ with
\[
\layer_n = \{\omega \in \Omega : \tau(\omega) \leq n\}, \quad \minLayer(\omega) = \tau(\omega)
\]
is a structure tower.
\end{proposition}

\noindent\textbf{Proof sketch.}
Covering: $\tau(\omega) \in \NN$ so $\omega \in \layer_{\tau(\omega)}$.
Monotonicity: $\tau(\omega) \leq i \leq j \Rightarrow \omega \in \layer_j$.
Minimality: $\tau(\omega)$ is the least $n$ with $\tau(\omega) \leq n$.
\qed

\begin{theorem}[Stopping sets are tower layers]
\label{thm:stopping-set-eq-layer}
The stopping set $\{\tau \leq n\}$ equals $\layer_n$ in $\Tower_\tau$.
\end{theorem}

\noindent\textbf{Proof.}
Immediate from Definition~\ref{def:tower-from-rank} and Proposition~\ref{prop:tower-from-stopping-time}.
\qed

\begin{theorem}[minLayer equals stopping time]
\label{thm:minlayer-eq-stopping}
For all $\omega \in \Omega$, $\minLayer_{\Tower_\tau}(\omega) = \tau(\omega)$.
\end{theorem}

\noindent\textbf{Proof.}
Apply Lemma~\ref{lem:minlayer-eq-rank} with $\rho = \tau$.
\qed

This result formalizes the intuition that ``the first time $\omega$ becomes observable'' (i.e., $\tau(\omega)$) is synonymous with ``the minimal layer containing $\omega$''.

\begin{example}[Constant stopping time]
\label{ex:constant-stopping}
Let $\tau(\omega) = K$ for all $\omega \in \Omega$.
Then $\layer_n = \emptyset$ for $n < K$, and $\layer_n = \Omega$ for $n \geq K$.
This is the ``step tower'' at level $K$.
\end{example}

In Lean (\texttt{StoppingTime\_C.lean}), the correspondence is made constructive via \texttt{towerFromStoppingTime}, with theorems \texttt{stoppingSet\_eq\_layer} and \texttt{minLayer\_eq\_stoppingTime} establishing the identification.

\subsection{Rank interpretation of filtrations}

The filtration $\Filtration$ itself induces a tower of $\sigma$-algebras.
For each singleton $\{\omega\}$, define the \emph{first measurable time}:
\[
\mathrm{firstMeas}(\omega) = \min\{n \in \NN : \{\omega\} \in \Filtration_n\}.
\]
This yields a rank function on $\Omega$, and the corresponding tower's layers are precisely the sets of singletons measurable at time $n$.

\section{Rank Optional Stopping Theorem}
\label{sec:rank-ost}

We now state and prove the Optional Stopping Theorem in the language of rank functions and structure towers.
The key insight is that the stopped process $M^\tau$ can be understood as ``restricting $M$ to layers of $\Tower_\tau$'', and the martingale property is preserved within layers.

\subsection{Stopped processes in rank language}

Recall the stopped process (Definition~\ref{def:stopped-process}):
\[
M^\tau_n(\omega) = M_{\min(n, \tau(\omega))}(\omega).
\]
In rank terms, $M^\tau_n(\omega)$ evaluates $M$ at the layer intersection: if $\omega \in \layer_n^{\Tower_\tau}$ and $\tau(\omega) \leq n$, then $M^\tau_n(\omega) = M_{\tau(\omega)}(\omega)$; otherwise, $M^\tau_n(\omega) = M_n(\omega)$.

\begin{lemma}[Stopped process is adapted]
\label{lem:stopped-adapted}
If $\tau$ is a stopping time and $M$ is adapted, then $M^\tau$ is adapted to $\Filtration$.
\end{lemma}

\noindent\textbf{Proof sketch.}
For each $n$, $M^\tau_n$ is measurable with respect to $\Filtration_n$ because:
\[
M^\tau_n = \sum_{k=0}^n M_k \cdot \mathbf{1}_{\{\tau = k\}} + M_n \cdot \mathbf{1}_{\{\tau > n\}}.
\]
Each indicator is $\Filtration_n$-measurable by the stopping-time property.
\qed

\begin{lemma}[Stopped process is integrable]
\label{lem:stopped-integrable}
If $M$ is integrable and $\tau$ is bounded by $K$, then $M^\tau$ is integrable.
\end{lemma}

\noindent\textbf{Proof sketch.}
$|M^\tau_n| \leq \max_{k \leq K} |M_k|$, which is integrable.
\qed

\subsection{The rank Optional Stopping Theorem}

The bounded OST asserts that stopping preserves expectation.

\begin{theorem}[Bounded Optional Stopping Theorem, rank version]
\label{thm:rank-ost-bounded}
Let $M$ be a martingale and $\tau$ a stopping time bounded by $K \in \NN$.
Then for all $n \in \NN$,
\[
\EE[M^\tau_n] = \EE[M_0].
\]
In particular, $\EE[M_{\tau}] = \EE[M_0]$.
\end{theorem}

\noindent\textbf{Proof sketch.}
\textbf{Step 1 (Stopped martingale):} $M^\tau$ is a martingale under the bounded-stopping hypothesis.
For the martingale property, compute:
\[
\EE[M^\tau_{n+1} \mid \Filtration_n] = \EE[M_{\min(n+1, \tau)} \mid \Filtration_n].
\]
Split into $\{\tau \leq n\}$ (where $M^\tau_{n+1} = M_\tau = M^\tau_n$) and $\{\tau > n\}$ (where $M^\tau_{n+1} = M_{n+1}$ and $\EE[M_{n+1} \mid \Filtration_n] = M_n = M^\tau_n$).

\textbf{Step 2 (Martingale expectation constancy):} For any martingale $H$, $\EE[H_n]$ is constant in $n$ by the tower property of conditional expectation:
\[
\EE[H_{n+1}] = \EE[\EE[H_{n+1} \mid \Filtration_n]] = \EE[H_n].
\]

\textbf{Step 3 (Conclusion):} Applying Step~2 to $H = M^\tau$,
\[
\EE[M^\tau_n] = \EE[M^\tau_0] = \EE[M_0].
\]
For the stopping-time expectation, note that if $\tau \leq K$, then $M^\tau_K(\omega) = M_{\tau(\omega)}(\omega)$ for all $\omega$, so $\EE[M_\tau] = \EE[M^\tau_K] = \EE[M_0]$.
\qed

\begin{remark}[Layer interpretation]
\label{rem:layer-interp}
In tower language, Theorem~\ref{thm:rank-ost-bounded} says: integrating over layers of $\Tower_\tau$ preserves the martingale's expected value.
Each layer $\layer_n = \{\tau \leq n\}$ is a ``stopping shell'', and the expectation is constant across shells.
\end{remark}

\begin{theorem}[Expectation at stopping time]
\label{thm:ost-at-tau}
Under the hypotheses of Theorem~\ref{thm:rank-ost-bounded},
\[
\EE\left[\sum_{\omega \in \Omega} M_{\tau(\omega)}(\omega) \cdot \mu(\{\omega\})\right] = \EE[M_0].
\]
In the discrete, finite setting, this is equivalent to $\EE[M_\tau] = \EE[M_0]$.
\end{theorem}

\begin{corollary}[Time-invariance of stopped expectation]
\label{cor:time-invariance}
For bounded $\tau$, the family $\{\EE[M^\tau_n]\}_{n \in \NN}$ is constant.
\end{corollary}

\subsection{Rank-theoretic perspective}

The bounded OST can be rephrased:
\begin{quote}
\emph{A rank function $\tau$ on $\Omega$ (satisfying stopping-time measurability) defines a tower $\Tower_\tau$.
The stopped process $M^\tau$ evaluates $M$ at the minimal layer for each $\omega$.
The martingale property ensures that layer-wise integration yields the same value as integrating over the initial layer.}
\end{quote}

This perspective suggests generalizations: any ``layer-respecting'' operation on martingales should preserve expectation, a principle exploited in advanced martingale transforms.

\section{Lean Mechanization and Executability}
\label{sec:lean}

The entire theory is formalized in Lean~4 with Mathlib.
The formalization consists of four main files, corresponding to Sections~\ref{sec:tower}--\ref{sec:rank-ost}.

\subsection{File structure}

\begin{enumerate}
\item \textbf{RankTower.lean} (Section~\ref{sec:tower}):
  \begin{itemize}
  \item Defines \texttt{StructureTowerWithMin} with fields: \texttt{carrier}, \texttt{layer}, \texttt{covering}, \texttt{monotone}, \texttt{minLayer}, \texttt{minLayer\_mem}, \texttt{minLayer\_minimal}.
  \item Constructs \texttt{towerFromRank} and \texttt{rankFromTower}.
  \item Proves the rank--tower isomorphism (Theorem~\ref{thm:rank-tower-iso}).
  \item Verifies the Vec2Q example (Example~\ref{ex:vec2q}).
  \end{itemize}

\item \textbf{StoppingTime\_C.lean} (Section~\ref{sec:stopping-rank}):
  \begin{itemize}
  \item Defines \texttt{stoppingTimeToRank} and \texttt{towerFromStoppingTime}.
  \item Proves \texttt{stoppingSet\_eq\_layer} (Theorem~\ref{thm:stopping-set-eq-layer}).
  \item Proves \texttt{minLayer\_eq\_stoppingTime} (Theorem~\ref{thm:minlayer-eq-stopping}).
  \item Implements constant stopping times (Example~\ref{ex:constant-stopping}).
  \end{itemize}

\item \textbf{Martingale\_RankStructure.lean} (martingale wrappers):
  \begin{itemize}
  \item Defines \texttt{rankStoppedProcess} as a thin wrapper around the existing stopped-process definition.
  \item Proves \texttt{rankOptionalStopping\_bounded}: the bounded OST in rank language.
  \item Proves adaptedness, integrability, and martingale-property preservation.
  \end{itemize}
  This file serves as a ``translation layer'', re-exporting existing Mathlib/project lemmas in rank-theoretic notation.

\item \textbf{OptionalStopping\_RankTheory.lean} (Section~\ref{sec:rank-ost}):
  \begin{itemize}
  \item States \texttt{rankOST\_bounded} (Theorem~\ref{thm:rank-ost-bounded}).
  \item States \texttt{rankOST\_at\_stopping\_time} (Theorem~\ref{thm:ost-at-tau}).
  \item States \texttt{rankOST\_expectation\_constant} (Corollary~\ref{cor:time-invariance}).
  \item Provides a layer-separated interface: depends only on \texttt{Martingale\_RankStructure.lean}, not on lower-level implementation files.
  \end{itemize}
\end{enumerate}

\subsection{Executability and decision procedures}

In the finite, rational-probability setting, all definitions are computable (modulo \texttt{Classical.choice} for $\minLayer$).
Lean's \texttt{\#eval} mechanism can verify small examples.
For instance, given a finite $\Omega = \{1,2,3\}$ with a two-step filtration and a simple martingale, one can compute:
\begin{itemize}
\item $\minLayer(\omega)$ for each $\omega$,
\item $M^\tau_n(\omega)$ for each $n, \omega$,
\item $\EE[M^\tau_n]$ by summing over $\Omega$.
\end{itemize}

This yields a \emph{certified calculator} for discrete martingale stopping problems.

\subsection{Proof engineering and layer separation}

A key design principle is \emph{layer separation}.
The files are organized so that high-level theorems (e.g., \texttt{rankOST\_bounded}) depend on \emph{thin wrapper} lemmas (e.g., \texttt{rankStopped\_martingale\_property}) rather than on low-level implementation details (e.g., sigma-algebra towers, conditional-expectation internals).
This makes the formalization modular and maintainable.

For example, \texttt{OptionalStopping\_RankTheory.lean} imports \texttt{Martingale\_RankStructure.lean} but not \texttt{StoppingTime\_MinLayer.md} or \texttt{Martingale\_StructureTower.md}.
The proof of \texttt{rankOST\_bounded} is essentially a one-liner calling the wrapper lemma that encapsulates the full proof.

This strategy minimizes proof obligations in the top-level file and clarifies the conceptual dependency graph.

\subsection{Limitations and future work}

\begin{itemize}
\item \textbf{Unbounded stopping times:} The formalization currently covers only bounded stopping times.
  Extending to unbounded $\tau$ requires uniform integrability and limit arguments, which are technically more involved.
\item \textbf{Continuous time:} Our setting is discrete.
  Extending rank towers to continuous time (Brownian filtrations, etc.) would require handling uncountable index sets and more sophisticated measure theory.
\item \textbf{Other martingale theorems:} Doob's maximal inequality, martingale convergence, and Burkholder--Davis--Gundy could be formulated in rank language.
\item \textbf{Category theory:} A full categorical treatment (Section~\ref{app:cat}) would formalize the rank--tower correspondence as an equivalence of categories, clarifying the ``naturality'' of our constructions.
\end{itemize}

\section{Discussion and Future Directions}
\label{sec:discussion}

We have demonstrated that the Optional Stopping Theorem admits a purely rank-theoretic formulation, where stopping times are rank functions and stopping sets are tower layers.
This perspective offers several advantages:

\begin{itemize}
\item \textbf{Conceptual clarity:} The layer-based view makes explicit the ``stratification'' of the sample space by stopping times.
\item \textbf{Algebraic simplicity:} Rank functions are elementary objects (maps to $\NN$), avoiding measure-theoretic subtleties in the discrete setting.
\item \textbf{Compositional proofs:} Layer separation allows modular, reusable proof components.
\item \textbf{Computability:} Finite, rational-probability examples are directly executable in Lean.
\end{itemize}

\paragraph{Connections to Bourbaki's mother structures.}
The rank--tower correspondence echoes Bourbaki's program of unifying mathematics via foundational structures (algebraic, order, topological).
Here, the structure tower serves as a ``mother structure'' for probability: filtrations, stopping times, and martingales all find natural expression as tower-theoretic objects.

\paragraph{Applications.}
Beyond OST, rank-based reasoning could illuminate:
\begin{itemize}
\item \textbf{Martingale transforms:} Predictable processes as rank-respecting operations.
\item \textbf{Doob decompositions:} Splitting processes into martingale and predictable parts via layer projections.
\item \textbf{Stochastic integration:} Discrete Itô calculus as a tower-theoretic construction.
\end{itemize}

\paragraph{Philosophical note.}
Formal verification forces precision.
By mechanizing the OST in Lean, we must specify every hypothesis, every definition, every lemma.
This discipline uncovers hidden assumptions in textbook proofs and ensures correctness.
The rank-based approach further simplifies verification by reducing probabilistic reasoning to combinatorial layer counting.

\section*{Acknowledgments}

This work was generated collaboratively, with \LaTeX{} typesetting by Claude Code (Anthropic) and Lean formalization by Codex (OpenAI).
The rank-tower framework was inspired by Bourbaki's structural approach and classical probability texts.

\begin{thebibliography}{99}

\bibitem{williams91}
D.~Williams.
\textit{Probability with Martingales}.
Cambridge University Press, 1991.

\bibitem{rogers-williams}
L.~C.~G.~Rogers and D.~Williams.
\textit{Diffusions, Markov Processes, and Martingales}, Volume 1.
Cambridge University Press, 2000.

\bibitem{kallenberg}
O.~Kallenberg.
\textit{Foundations of Modern Probability}, 2nd edition.
Springer, 2002.

\bibitem{doob}
J.~L.~Doob.
\textit{Stochastic Processes}.
Wiley, 1953.

\bibitem{isabelle-ost}
J.~Hölzl, A.~Lochbihler, and D.~Traytel.
A formalized theory of martingales in Isabelle/HOL.
\textit{Archive of Formal Proofs}, 2015.

\bibitem{coq-prob}
A.~Affeldt, R.~Affeldt, and D.~Nowak.
A Coq formalization of measure theory and probability.
\textit{Journal of Formalized Reasoning}, 2018.

\bibitem{lean-mathlib}
The mathlib Community.
The Lean mathematical library.
\url{https://github.com/leanprover-community/mathlib4}, 2024.

\bibitem{bourbaki}
N.~Bourbaki.
\textit{Éléments de Mathématique}.
Springer, various volumes, 1939--present.

\end{thebibliography}

\newpage
\appendix

\section{RankTower Normal Form}
\label{app:normal-form}

This appendix elaborates on the rank--tower correspondence (Theorem~\ref{thm:rank-tower-iso}), providing additional examples and a category-theoretic perspective.

\subsection{Normal form theorem (detailed)}

Every structure tower $\Tower$ admits a \emph{canonical rank function} $\rank_\Tower = \minLayer$, and reconstructing the tower from this rank yields the original tower (up to isomorphism).

\begin{theorem}[Normal form]
\label{thm:normal-form-detailed}
Let $\Tower$ be a structure tower.
Then $\Tower \cong \Tower_{\rank_\Tower}$ via the identity map on carriers, with layer correspondence $\layer_n = \{x : \rank_\Tower(x) \leq n\}$.
\end{theorem}

\noindent\textbf{Proof.}
Define $\rank_\Tower(x) = \minLayer(x)$.
By Definition~\ref{def:tower}, $x \in \layer_n$ iff $\minLayer(x) \leq n$, so $\layer_n = \{x : \rank_\Tower(x) \leq n\}$.
This is precisely $\layer_n^{\Tower_{\rank_\Tower}}$ by Definition~\ref{def:tower-from-rank}.
\qed

\subsection{Additional examples}

\begin{example}[Natural number ordering]
Let $X = \NN$, $\rho(n) = n$.
Then $\layer_k = \{0,1,\ldots,k\}$ and $\minLayer(n) = n$.
This is the ``discrete tower'' over $\NN$.
\end{example}

\begin{example}[Cardinality rank for finite sets]
Let $X = \mathcal{P}(\{1,2,\ldots,N\})$ (the power set).
Define $\rho(S) = |S|$.
Then $\layer_k$ is the family of subsets with at most $k$ elements.
This tower is used in combinatorics (e.g., the inclusion--exclusion principle can be phrased as a layer-by-layer computation).
\end{example}

\section{Category Theory: Lanes and Forgetful Functors}
\label{app:cat}

We sketch a categorical treatment, deferred from the main text to keep the exposition accessible.

\subsection{The category $\mathsf{Tower}$}

Define the category $\mathsf{Tower}$:
\begin{itemize}
\item \textbf{Objects:} Structure towers $\Tower = (X, \layer, \minLayer)$.
\item \textbf{Morphisms:} A morphism $f: \Tower_1 \to \Tower_2$ is a map $f: X_1 \to X_2$ such that:
  \begin{enumerate}
  \item $f(\layer_n^{\Tower_1}) \subseteq \layer_n^{\Tower_2}$ for all $n$ (layer preservation),
  \item $\minLayer_{\Tower_2}(f(x)) \leq \minLayer_{\Tower_1}(x)$ for all $x$ (rank non-increase).
  \end{enumerate}
\end{itemize}

\subsection{The category $\mathsf{Rank}$}

Define $\mathsf{Rank}$:
\begin{itemize}
\item \textbf{Objects:} Pairs $(X, \rho)$ where $\rho: X \to \NN$.
\item \textbf{Morphisms:} Maps $f: X_1 \to X_2$ with $\rho_2(f(x)) \leq \rho_1(x)$ (rank non-increase).
\end{itemize}

\subsection{Equivalence of categories}

The constructions $\Tower_\rho$ (Definition~\ref{def:tower-from-rank}) and $\rank_\Tower$ (Definition~\ref{def:rank-from-tower}) extend to functors:
\[
F: \mathsf{Rank} \to \mathsf{Tower}, \quad G: \mathsf{Tower} \to \mathsf{Rank}.
\]

\begin{proposition}
\label{prop:equivalence}
$F$ and $G$ form an equivalence of categories: $G \circ F \cong \mathrm{id}_{\mathsf{Rank}}$ and $F \circ G \cong \mathrm{id}_{\mathsf{Tower}}$.
\end{proposition}

\noindent\textbf{Proof sketch.}
Natural isomorphisms are given by Theorem~\ref{thm:rank-tower-iso}.
The rank-non-increase condition on morphisms ensures functoriality.
\qed

\subsection{Filtration as a tower functor}

A filtration $\Filtration = (\Filtration_n)_{n \in \NN}$ induces a functor from the poset category $(\NN, \leq)$ to the category of $\sigma$-algebras.
This can be lifted to a tower on $\Omega$ by associating to each $\omega$ its first-measurable-time rank.
The functoriality of $\Filtration$ ensures that rank functions arising from stopping times are coherent with the filtration's layer structure.

\subsection{Adjunctions and free constructions}

One can consider the forgetful functor $U: \mathsf{Tower} \to \mathsf{Set}$ (forgetting layers).
A left adjoint to $U$ would be the ``free tower'' construction, assigning to each set $X$ the trivial tower with $\layer_n = X$ for all $n$.
In probability, this corresponds to the trivial filtration (all sets measurable at time 0).

The rank--tower correspondence provides an alternative ``canonical'' free construction: given a rank function, freely generate the tower.
This suggests a categorical interpretation of the Optional Stopping Theorem as a naturality property of certain functors between tower-enriched categories.

\end{document}
